\newcommand{\tipoRequerimento}{Contestação do memorial de cálculo, pagamento complementar do dobro}

\section{DOS FATOS E DOS DIREITOS}

\begin{enumerate}[resume]
\item Conforme disposto no item 35 da Nota Técnica nº 136/2024, nos casos de condutas continuadas da distribuidora que resultem em faturamento incorreto e que abranjam períodos regidos pelas Resoluções Normativas nº 414/2010 e nº 1.000/2021, a análise acerca da devolução em dobro deve ser realizada de forma unificada. Ainda segundo o referido entendimento, quando a constatação da irregularidade ocorrer na vigência da REN nº 1.000/2021, a avaliação da hipótese de engano justificável prevista no art. 113 da REN nº 414/2010 deve observar as disposições estabelecidas no § 3º do art. 323 da REN nº 1.000/2021.
\item No caso em análise, conforme destacado no item 36 da Nota Técnica nº 136/2024, a constatação do faturamento a maior ocorreu durante a vigência da REN nº 1.000/2021. Ademais, o erro foi causado exclusivamente pela distribuidora, inexistindo qualquer elemento que permita atribuir responsabilidade ao consumidor ou a terceiros. Dessa forma, não se configura a hipótese de engano justificável prevista no § 3º do art. 323 da REN nº 1.000/2021. 
\item Embora a concessionária tenha efetuado o pagamento da quantia de  R\$ 1.889,39, apurado no memorial de cálculo, o mesmo se mostra incompleto, pois não incluiu a devolução em dobro dos últimos 5 (cinco) anos, em desacordo com a norma. Portanto, é essencial que a Ouvidoria da ANEEL observe e aplique os precedentes já estabelecidos pela própria Agência para garantir a quitação do valor complementar devido.
\end{enumerate}

\section{DOS PEDIDOS}

\begin{enumerate}[resume]
\item Ante o exposto, requer-se:
    \begin{enumerate}
    \item A observância e a aplicação dos precedentes administrativos já estabelecidos pela ANEEL, reconhecendo-se a incidência da devolução em dobro dos valores cobrados indevidamente, nos termos do § 3º do art. 323 da REN nº 1.000/2021 e da Nota Técnica nº 136/2024;
    \item A apuração e o pagamento da complementação dos valores devidos, considerando a devolução em dobro dos montantes cobrados indevidamente nos últimos 5 (cinco) anos, com a correspondente atualização monetária e demais acréscimos regulamentares cabíveis;

    \end{enumerate}
\end{enumerate}